% 2013 年安徽卷文科第 17 题：甲、乙两校各 30 名学生数学成绩的茎叶图。
% 来源 PDF：普通高考/2013/2013安徽文.pdf，第 1 页（印面 785）右栏。
% 茎为十位数字 4—9；甲校叶置左、乙校叶置右，逐位照录原图（未重新排序）。
\begin{tikzpicture}[font=\small, x=1cm, y=1cm]
  % 表头
  \node at (-0.7,0.52) {甲};
  \node at (2.35,0.52) {乙};
  \draw (-2.5,0.3) -- (3.9,0.3);
  % 双竖线与茎
  \draw (0.28,0.22) -- (0.28,-2.42);
  \draw (0.98,0.22) -- (0.98,-2.42);
  % 各行：茎 4—9
  \foreach \stem/\y in {4/0, 5/-0.44, 6/-0.88, 7/-1.32, 8/-1.76, 9/-2.2} {
    \node at (0.63,\y) {$\stem$};
  }
  % 甲校叶（右对齐）
  \node[anchor=east] at (0.17,0)    {$7$};
  \node[anchor=east] at (0.17,-0.44){$5\ 3\ 3\ 2$};
  \node[anchor=east] at (0.17,-0.88){$5\ 5\ 4\ 3\ 3\ 3\ 1\ 0\ 0$};
  \node[anchor=east] at (0.17,-1.32){$8\ 6\ 6\ 2\ 2\ 1\ 1\ 0\ 0$};
  \node[anchor=east] at (0.17,-1.76){$7\ 5\ 4\ 4\ 2$};
  \node[anchor=east] at (0.17,-2.2) {$2\ 0$};
  % 乙校叶（左对齐）
  \node[anchor=west] at (1.09,0)    {$5$};
  \node[anchor=west] at (1.09,-0.44){$3\ 3\ 8$};
  \node[anchor=west] at (1.09,-0.88){$0\ 0\ 0\ 1\ 1\ 2\ 2\ 3\ 3\ 5$};
  \node[anchor=west] at (1.09,-1.32){$0\ 0\ 2\ 2\ 2\ 3\ 3\ 6\ 6\ 9$};
  \node[anchor=west] at (1.09,-1.76){$1\ 1\ 5\ 5\ 8$};
  \node[anchor=west] at (1.09,-2.2) {$0$};
\end{tikzpicture}
